\documentclass{article}
\usepackage{graphicx}
\usepackage{amssymb}
\usepackage[colorlinks=true, linkcolor=blue]{hyperref}
\title{Mini Game Hub Project Report}
\author{Rasagnya}
\date{\today}
\begin{document}
\maketitle
\tableofcontents
\newpage

\section{Introduction}
This is a project made as part of the CS108 course (Spring semester 2026) for gaming purpose.
This Mini Game Hub project includes three classic board games:

\begin{itemize}
    \item It is a simple gaming application that combines three games: Tic Tac Toe, Connect 4, and Othello.
    
    \item The main aim is to provide a single platform where users can easily play different games.
    
    \item Each game helps in improving thinking and strategy skills in a fun way.
    
    \item Tic Tac Toe is a basic game that helps understand simple rules and patterns.
    
    \item Connect 4 is a game where players need to plan their moves and think ahead.
    
    \item Othello is a more advanced game that involves capturing opponent pieces using strategy.
    
    \item The project focuses on making the games easy to use with simple controls and clear rules.
    
    \item It also helps in learning how different game logics and states are implemented in programming.
\end{itemize}
\section{Running Instructions}
To run this Mini Game Hub project, follow these steps:
\begin{itemize}
    \item Make sure python 3.10 and git bash are available in your computer.
    \item  Clone the project repository or download the source code, and open the project root directory in a terminal.
    \item Install the required libraries like Numpy , Matplotlib etc.
    \item Run this application using the entry-point script main.sh:
    bash main.sh.
\end{itemize}
\section{Game Hub Layout}
\begin{itemize}
    \item The Game Hub starts with a main menu screen where users can choose between different games.

    \item The available game options are Tic Tac Toe, Connect 4, and Othello, displayed as clickable buttons or selections.

    \item After selecting a game, the user is taken to a separate game screen where that particular game is played.

    \item Each game has its own dedicated layout, rules, and interaction system within the same project.

    \item The main menu acts as a central navigation point for switching between games.

    \item The overall design is kept simple and clean so that users can easily understand and navigate between different games.

    \item A consistent UI style is maintained across all games for better user experience and smooth interaction.
\end{itemize}
\section{Modules Used}
\cite{iitb_latex}
\subsection{Numpy}

\subsection{Pygame}
\subsection{sys}
\subsection{Matplotlib}
\subsection{Time}
\subsection{os}
\subsection{subprocess}
\begin{itemize}

\item \textbf{NumPy:} NumPy\cite{numpy_docs} is a Python library used for numerical computing and working with arrays. It helps in performing fast mathematical operations on large data sets.

\item \textbf{Pygame:} Pygame\cite{pygame_docs} is a Python library used for creating games and multimedia applications. It provides tools for graphics, sound, and handling user input.

\item \textbf{sys:} The sys module provides access to system-specific functions and parameters in Python. It is commonly used for command-line arguments and controlling program execution.

\item \textbf{Matplotlib:} Matplotlib\cite{matplotlib_docs} is a Python library used for data visualization and plotting graphs. It helps in creating charts like line plots, bar graphs, and histograms.

\item \textbf{Time:} The time module is used to handle time-related functions in Python. It is useful for adding delays, measuring execution time, and working with timestamps.

\item \textbf{os:} The os module provides functions to interact with the operating system. It is used for file handling, directory operations, and running system commands.

\item \textbf{subprocess:} The subprocess module is used to run external programs and system commands from Python. It helps in integrating Python with other applications or scripts.

\end{itemize}
\section{Features Implemented}
\begin{itemize}
\item{Dynamic grid size implemented for all 3 games Tictactoe,Othello,Connect4\\  i.e we can change size of grid to any $n \times m$ where $m,n \in \mathbb{N}$}
\item{Disc Movement Feature in Connect4:\\A preview disc appears at top ,when the player clicks,the disc is dropped into the lowest available position}
\item{Othello Highlighting:\\ In OPthello if player1 is playing the game his side highlights avoiding confusion.}
\item{Showing the statistics,top5 player and their win count using bar graph and games and their frequency using pie chart by using matplotlib}
\end{itemize}
\section{Code Logic}
\subsection{Tic-Tac-Toe}
The game runs in a loop where players take turns clicking on the grid to place their symbol. Each click fills an empty box, and after every move, the game checks if the player has formed a line of five symbols in any direction or if the board is completely filled. If someone wins or it’s a draw, the result is shown and saved, and the game ends. If not, the turn switches to the other player and the game continues.
\subsection{Othello}
The Othello game logic initializes an 8×8 board and updates it based on valid player moves using directional checks for flipping opponent pieces. After every move, the board is checked for game termination either when the board is full or when no valid moves exist for both players. When the game ends, the piece counts are compared to determine the winner or a draw. The result is then displayed, stored in history, and used to update the leaderboard and visualization.
\subsection{Connect4}
The game keeps running in a loop, watching for mouse movement and clicks. When a player clicks on a column, the program places their piece in the lowest empty spot of that column. After placing the piece, it checks if that move created a line of four pieces or if the board is full.
If the game is over, it decides the winner (or draw), saves the result, shows it on the screen, and stops the game. If not, it switches to the other player and continues the loop.
\subsection{Authentication System }
main.sh handles user authentication before starting the game. It prompts each player to either log in or register, storing passwords securely as hashed values in a users.tsv file and verifying them during login. If a username doesn’t exist, it offers registration; if the password is wrong, it keeps asking until correct. Finally, it ensures both players are different users and then runs the Python game with their usernames as arguments \cite{bash_docs}.
\subsection{Leaderboard}
leaderboard.sh reads game results from history.csv and builds a leaderboard by counting wins and losses for each user in each game. It skips draws, calculates a win/loss ratio (or inf if no losses), and prints a formatted table for games like TicTacToe, Othello, and Connect4. The results are then sorted dynamically based on a chosen metric (wins, losses, or ratio) passed as an argument. Finally, it displays the sorted leaderboard along with the metric used.
\section{Hurdles Faced}
\begin{itemize}
\item{Since the code for any file is very large we have faced huge problems with syntax . Also ideas for games are difficult and git push and pull part is the most difficult part , because many times we lost whatever code we have written in this git push part.}
\item{Got errors when running leaderboard.sh inside games.got Insufficient buffer space queue full error,the error is caused due to forcing python to use git bash from inside program,We solved it by running the program from CMD}
\item{Circular Import.We were importing game.py into tictactoe,connect4,othello from and tictactoe,othello,connect4 into game.py.\\We solved this by defining the baseclass in another file "baseclass.py" }
\item Encountered multiple syntax and indexing errors. For example, using \texttt{win\_name = players[winner]} caused issues because \texttt{winner} was a string value ("1" or "-1"). While "1" worked correctly, "-1" led to unintended negative indexing. This was resolved by changing the approach to \texttt{win\_name = self.getCurrentPlayer()}.

\item While reading from \texttt{history.csv}, the values \texttt{winner, loser, date, game} contained a trailing newline character (\texttt{\textbackslash n}), which caused errors during processing. This issue was fixed by applying the \texttt{.strip()} function before splitting the line.
\end{itemize}
\section{Future Improvements}
We will try to add one more game in the mini game hub . We will try to improve the design of our game hub so that it gives you gaming vibes . Make it more generalised . We will try to include some animations wherever required.
\bibliographystyle{plain}
\bibliography{ref}
\end{document}
